\documentclass[a4paper,11pt]{article}
\usepackage[utf8]{inputenc}
\usepackage[T1]{fontenc}
\usepackage{tgheros}
\renewcommand{\familydefault}{\sfdefault}
\usepackage[colorlinks=true, urlcolor=blue]{hyperref}
\usepackage{geometry}
\geometry{left=1in, right=1in, top=1in, bottom=1in}
\usepackage{setspace}
\onehalfspacing

\begin{document}

\begin{center}
\textbf{\large How Equal is Equal Opportunity?} \\[4pt]
\textbf{Gender, Ethnic and Regional Disparities in Labor Market Outcomes \\
of Higher Education Policies in Brazil} \\[10pt]
\textit{Baha Khammari} \\
M.Sc.\ Economics, University of Cologne \\
Supervisors: Prof.\ Dr.\ Pia Pinger \& Anna Person, M.Sc. \\
Grade: 1.0/1.0 \\[6pt]
April 2026
\end{center}

\vspace{12pt}

\noindent
This thesis estimates the causal labour-market effects of Brazil's ProUni scholarship programme on formal-sector wages across 558 microregions over the period 2005--2019. I apply the Callaway \& Sant'Anna (2021) heterogeneity-robust staggered difference-in-differences estimator with doubly-robust inference to a microregion-level panel constructed by linking ProUni administrative records, RAIS formal employer--employee data ($\sim$5\,million observations), and Census demographics via the \texttt{basedosdados} BigQuery interface.

The analysis documents a non-monotonic dose--response relationship: Low-Dose regions exhibit a statistically significant +3.04\,\% wage effect (DR estimate, $p < 0.10$), while Mid-Dose regions show negative effects coinciding with recession timing, and High-Dose regions display near-zero effects consistent with diminishing returns. A na\"ive two-way fixed-effects (TWFE) specification masks these heterogeneous effects entirely ($\hat{\beta}_{\text{TWFE}} = -0.65\%$), providing an empirical illustration of the ``forbidden comparisons'' bias identified by Goodman-Bacon (2021).

Heterogeneity analysis reveals that wage gains accrue disproportionately to non-white workers in moderate-exposure markets, consistent with ProUni's affirmative-action design, while a gender asymmetry persists despite majority-female scholarship uptake. Robustness checks include HonestDiD sensitivity analysis (Rambachan \& Roth, 2023), age-cohort placebo tests, Bacon decomposition, alternative dose cutoffs, and concurrent-policy controls (FIES, Quota Law).

\vspace{8pt}
\noindent\textbf{Keywords:} ProUni, difference-in-differences, staggered adoption, doubly-robust estimation, labour markets, Brazil, higher education policy.

\vspace{4pt}
\noindent\textbf{JEL Codes:} I26, J31, C23, O15.

\end{document}
